Likewise, one lets \(G\) act on the right on
\(\End_{I_{S}(\mathcal{M})/S}(G_{I_{S}(\mathcal{M})})\) as follows: for every
\(S'\to S\), \(x\in G(S')\) and
\(u\in\End_{I_{S}(\mathcal{M})/S}(G_{I_{S}(\mathcal{M})})(S')
=\End_{I_{S'}(\mathcal{M})}(G_{I_{S'}(\mathcal{M})})\),
\[
(u\cdot x)(g)=u(g\cdot x^{-1})\cdot x,
\]
for every \(g\in G(S'')\), \(S''\to I_{S'}(\mathcal{M})\). Then the morphism
of~3.12.1
\[
\Hom_{G/S}\bigl(G,T_{G/S}(\mathcal{M})\bigr)
\longrightarrow
\End_{I_{S}(\mathcal{M})/S}\bigl(G_{I_{S}(\mathcal{M})}\bigr)
\]
respects the operations of \(G\) and therefore induces, for every \(S'\to S\),
a bijection between \(\Gamma(T_{G_{S'}/S'}(\mathcal{M})/G_{S'})\) and the set of
\(I_{S'}(\mathcal{M})\)-endomorphisms \(u\) of \(G_{I_{S'}(\mathcal{M})}\) which
induce the identity on \(G\) and which ``commute with the right
translations'', \ie which satisfy \(u_{S''}\cdot x=u_{S''}\) for every
\(S''\to S'\) and \(x\in G(S'')\). One therefore obtains:

\begin{proposition}{4.1.3}\label{II.4.1.3}
\((*)\) There exists a bijection, functorial in \(G\), between the set
\[
\Lie(G/S,\mathcal{M})(S)
\]
and the set of
\(I_{S}(\mathcal{M})\)-endomorphisms of \(G_{I_{S}(\mathcal{M})}\) inducing the
identity on \(G\) and commuting with the right translations of \(G\) (in the
sense indicated above).
\end{proposition}

Now taking~3.13 into account:

\begin{theorem}{4.1.4}\label{II.4.1.4}
\((*)\) Let \(G\) be an \(S\)-group functor; suppose that \(G/S\) satisfies~(E).
Then the group \(\Lie(G/S,\mathcal{M})(S)\) is identified, functorially in
\(G\), with the group of \(I_{S}(\mathcal{M})\)-automorphisms of
\(G_{I_{S}(\mathcal{M})}\) inducing the identity on \(G\) and commuting with
the right translations of \(G\) (in the sense indicated above).
\end{theorem}

One thus recovers (in the case \(\mathcal{M}=\mathbf{O}_{S}\)) one of the
classical definitions of the Lie algebra of a group.

\numpara{4.2.0}\label{II.4.2.0}%
{\renewcommand{\thefootnote}{(54)}\nde{One has added paragraph~4.2.0, whose
results are used implicitly in~4.2 and~4.7 of the original.}}
Before going further, let us establish new corollaries to~3.11. Let \(X\),
\(Y\) be over \(Z\), and \(Z\) over \(S\), as in section~1. As one saw
in~3.11, the isomorphisms~1.1~(4):
\[
(1)\qquad
\resizebox{0.9\linewidth}{!}{\(\vcenter{\xymatrix@C=3.5em@R=2.4em{
\Hom_{S}\bigl(I_{S}(\mathcal{M}),\Hom_{Z/S}(X,Y)\bigr)
  \ar[rr]^{\sim}
  \ar[dr]_{\simeq}
&&
\Hom_{Z/S}\bigl(X,\Hom_{Z}(I_{Z}(\mathcal{M}),Y)\bigr)
  \ar[dl]^{\simeq}
\\
&
\Hom_{Z/S}\bigl(X\times_{S}I_{S}(\mathcal{M}),Y\bigr)
}}\)}
\]
induce the isomorphism \(\theta\) below
\[
(2)\qquad
\resizebox{0.9\linewidth}{!}{\(\vcenter{\xymatrix@C=3.5em@R=2.4em{
T_{\Hom_{Z/S}(X,Y)}(\mathcal{M})
  \ar[rr]^{\theta}_{\sim}
  \ar[dr]_{\simeq}
&&
\Hom_{Z/S}\bigl(X,T_{Y/Z}(\mathcal{M})\bigr)
  \ar[dl]^{\simeq}
\\
&
\Hom_{Z/S}\bigl(X\times_{S}I_{S}(\mathcal{M}),Y\bigr)
.}}\)}
\]

By~1.3, if \(Y\) is a \(Z\)-group, the same is true of \(\Hom_{Z}(V,Y)\) for
every \(V\to Z\) (in particular for \(V=I_{Z}(\mathcal{M})\)); explicitly, if
\(Z''\to Z'\to Z\) and \(\varphi,\psi\in\Hom_{Z}(V_{Z'},Y)\), then
\(\varphi\cdot\psi\) is defined by \((\varphi\cdot\psi)(v)=\varphi(v)\psi(v)\),
for every \(v\in V_{Z'}(Z'')\).

Suppose now that \(X\) and \(Y\) are \(Z\)-groups, and let us set the following
definition.

\begin{definition}{4.2.0.1}\label{II.4.2.0.1}
Let \(\Hom_{(Z/S)\text{-gr.}}(X,Y)\) be the subfunctor of
\(\Hom_{Z/S}(X,Y)\) defined by: for every \(S'\to S\),
\begin{equation}
\tag{3}
\Hom_{(Z/S)\text{-gr.}}(X,Y)(S')
=
\Hom_{Z_{S'}\text{-gr.}}(X_{S'},Y_{S'}).
\end{equation}
This definition applies equally when one replaces \(Y\) by the \(Z\)-group
\(T_{Y/Z}(\mathcal{M})\).
\end{definition}

One then sees easily that
\(T_{\Hom_{(Z/S)\text{-gr.}}(X,Y)}(\mathcal{M})(S')\) corresponds, in the
preceding isomorphisms~(2), to the \(Z_{S'}\)-morphisms
\[
\varphi\colon X_{S'}\times_{S'}I_{S'}(\mathcal{M})\longrightarrow Y_{S'}
\]
which are ``multiplicative in \(X\)'', \ie which satisfy
\(\varphi(x_{1}x_{2},m)=\varphi(x_{1},m)\varphi(x_{2},m)\), and that these
correspond to the morphisms of \(Z_{S'}\)-groups
\(X_{S'}\to T_{Y/Z}(\mathcal{M})_{S'}\). One has thus obtained:

\begin{proposition}{4.2.0.2}\label{II.4.2.0.2}
Let \(X\), \(Y\) be \(Z\)-groups, with \(Z\) over \(S\). One has isomorphisms
of \(S\)-functors, functorial in \(\mathcal{M}\):
\[
T_{\Hom_{(Z/S)\text{-gr.}}(X,Y)}(\mathcal{M})
\isomto
\Hom_{(Z/S)\text{-gr.}}\bigl(X,T_{Y/Z}(\mathcal{M})\bigr).
\]
\end{proposition}

In particular, for \(Z=S\), one obtains the following corollary. Before
stating it, let us remark that if \(Y\) is a commutative \(S\)-group, the
same is true of \(T_{Y/S}(\mathcal{M})\), then of
\(H=\Hom_{S\text{-gr.}}(X,Y)\) and
\(\Hom_{S\text{-gr.}}(X,T_{Y/S}(\mathcal{M}))\), and finally of
\(T_{H/S}(\mathcal{M})\).

\begin{corollary}{4.2.0.3}\label{II.4.2.0.3}
Let \(X\), \(Y\) be \(S\)-groups. One has isomorphisms of \(S\)-functors,
functorial in \(\mathcal{M}\):
\[
T_{\Hom_{S\text{-gr.}}(X,Y)/S}(\mathcal{M})
\isomto
\Hom_{S\text{-gr.}}\bigl(X,T_{Y/S}(\mathcal{M})\bigr).
\]
If \(Y\) is commutative, these are moreover isomorphisms of abelian
\(S\)-groups.
\end{corollary}

\begin{definition}{4.2.0.4}\label{II.4.2.0.4}
{\renewcommand{\thefootnote}{(55)}\nde{One has inserted here this
definition, which in the original appeared in~4.3.}}
If \(Y\) is an \(\mathbf{O}_{S}\)-module, the functor
\(T_{Y/S}(\mathcal{M})\) (\resp \(\Lie(Y/S,\mathcal{M})\)) is equipped with an
\(\mathbf{O}_{S}\)-module structure deduced from that of \(Y\). Equipped
with this structure, one will denote it by \(T'_{Y/S}(\mathcal{M})\)
(\resp \(\Lie'(Y/S,\mathcal{M})\)).
\end{definition}

Consequently, if \(X\), \(Y\) are \(\mathbf{O}_{S}\)-modules, then
\(T'_{Y/S}(\mathcal{M})=\Hom_{S}(I_{S}(\mathcal{M}),Y)\) and
\(H=\Hom_{\mathbf{O}_{S}\text{-mod.}}(X,Y)\), and then
\(\Hom_{\mathbf{O}_{S}\text{-mod.}}(X,T'_{Y/S}(\mathcal{M}))\) and
\(T'_{H/S}(\mathcal{M})\), are equipped with an \(\mathbf{O}_{S}\)-module
structure, and one obtains the following:

\begin{corollary}{4.2.0.5}\label{II.4.2.0.5}
If \(X\), \(Y\) are \(\mathbf{O}_{S}\)-modules, one has isomorphisms of
\(\mathbf{O}_{S}\)-modules, functorial in \(\mathcal{M}\):
\[
T'_{\Hom_{\mathbf{O}_{S}\text{-mod.}}(X,Y)/S}(\mathcal{M})
\isomto
\Hom_{\mathbf{O}_{S}\text{-mod.}}\bigl(X,T'_{Y/S}(\mathcal{M})\bigr).
\]
\end{corollary}

\begin{definition}{4.2.A}\label{II.4.2.A}
{\renewcommand{\thefootnote}{(56)}\nde{One has expanded the original in what
follows; in particular, one has added definition~4.2.A and
remark~4.2.B.}}
Let \(X\), \(L\) be \(S\)-groups, with \(X\) operating on \(L\) by group
automorphisms (\cf I~2.3.5). One defines the subfunctor \(Z^{1}_{S}(X,L)\)
of \(\Hom_{S}(X,L)\) as follows: for every \(S'\to S\),
\[
Z^{1}_{S}(X,L)(S')
=
\left\{
\varphi\in\Hom_{S'}(X_{S'},L_{S'})
\;\middle|\;
\begin{aligned}
&\varphi(x_{1}x_{2})=\varphi(x_{1})\,(x_{1}\cdot\varphi(x_{2})),\\
&\text{for every }x_{1},x_{2}\in X(S''),\ S''\to S'
\end{aligned}
\right\}.
\]
One calls it the ``functor of crossed homomorphisms from \(X\) to \(L\)''.
\end{definition}

\begin{remark}{4.2.B}\label{II.4.2.B}
\textsuperscript{(56)}~a)~If \(L'\) is a second \(S\)-group on which \(X\) operates by group
automorphisms, one has
\[
Z^{1}_{S}(X,L\times_{S}L')
\simeq
Z^{1}_{S}(X,L)\times_{S}Z^{1}_{S}(X,L').
\]
b) If \(L\) is a \(G\)-\(\mathbf{O}_{S}\)-module, % typo?: kept as printed (G-O_S-module, not X)
\(Z^{1}_{S}(X,L)\) coincides with the kernel of the differential
\[
\partial\colon\Hom_{S}(X,L)\to\Hom_{S}(X^{2},L)
\]
defined in~I~5.1; in
particular, \(Z^{1}_{S}(X,L)\) is in this case an \(\mathbf{O}_{S}\)-module.
\end{remark}

\numpara{4.2}\label{II.4.2}
Let \(u\colon X\to Y\) be a morphism of \(S\)-groups; then
\(L^{u}_{\Hom_{S\text{-gr.}}(X,Y)/S}(\mathcal{M})\) is described as follows.
First, one saw in~3.11.3 that one has an isomorphism of \(S\)-functors,
functorial in \(\mathcal{M}\):
\[
(\dagger)\qquad
L^{u}_{\Hom_{S}(X,Y)/S}(\mathcal{M})
\isomto
\Hom_{Y/S}\bigl(X,T_{Y/S}(\mathcal{M})\bigr).
\]
On the other hand, as \(Y\) is an \(S\)-group; one then has
\[
T_{Y/S}(\mathcal{M})
=
\Lie(Y/S,\mathcal{M})\cdot Y
=
\Lie(Y/S,\mathcal{M})_{Y}\,;
\]
there results an isomorphism of \(S\)-functors, functorial in \(\mathcal{M}\):
\[
L^{u}_{\Hom_{S}(X,Y)/S}(\mathcal{M})
\isomto
\Hom_{S}\bigl(X,\Lie(Y/S,\mathcal{M})\bigr).
\]
For every \(S'\to S\), let us denote by \(u'\colon X'\to Y'\) the morphism
deduced from \(u\) by base change. Let us consider the \(S\)-functor defined
as follows:
{\renewcommand{\thefootnote}{(57)}\nde{This is definition~4.2.0.1 applied
to \(Z=Y\) and to the \(Y\)-groups \(u\colon X\to Y\) and
\(T_{Y/S}(\mathcal{M})\).}}
for every \(S'\to S\),
\begin{align*}
\Hom_{(Y/S)\text{-gr.}}\bigl(X,\Lie(Y/S,\mathcal{M})\cdot Y\bigr)(S')
&=
\Hom_{Y'\text{-gr.}}\bigl(X',(\Lie(Y/S,\mathcal{M})\cdot Y)_{S'}\bigr)\\
&=
\Hom_{Y'\text{-gr.}}\bigl(X',\Lie(Y'/S',\mathcal{M})\cdot Y'\bigr).
\end{align*}
Then one sees easily that the isomorphism \((\dagger)\) induces an
isomorphism
\[
(\dagger')\qquad
L^{u}_{\Hom_{S\text{-gr.}}(X,Y)/S}(\mathcal{M})
\isomto
\Hom_{(Y/S)\text{-gr.}}\bigl(X,\Lie(Y/S,\mathcal{M})\cdot Y\bigr).
\]
On the other hand, the morphism
\[
X\xrightarrow{u}Y\xrightarrow{\Ad}
\Aut_{S\text{-gr.}}\bigl(\Lie(Y/S,\mathcal{M})\bigr)
\]
defines an operation of \(X\) on \(L=\Lie(Y/S,\mathcal{M})\) by group
automorphisms.

If \(\Phi\in\Hom_{Y/S}\bigl(X,\Lie(Y/S,\mathcal{M})\cdot Y\bigr)\) then, for
every \(S''\to S'\to S\) and \(x\in X(S'')\), one can write in a unique way
\[
\Phi(S')(x)=\varphi(S')(x)\cdot u'(x),
\qquad\text{where}\qquad
\varphi(S')(x)\in\Lie(Y'/S',\mathcal{M})(S'');
\]
this determines an element \(\varphi\) of
\(\Hom_{S}\bigl(X,\Lie(Y/S,\mathcal{M})\bigr)\). Then \(\Phi(S')\) is a
morphism of groups if, and only if, for every \(x_{1},x_{2}\in X(S'')\) one
has:
\begin{align*}
\varphi(S')(x_{1}x_{2})
&=
\varphi(S')(x_{1})\,
\bigl(u(x_{1})\,\varphi(S')(x_{2})\,u(x_{1})^{-1}\bigr)\\
&=
\varphi(S')(x_{1})\,\bigl(x_{1}\cdot\varphi(S')(x_{2})\bigr),
\end{align*}
\ie if and only if \(\varphi\in Z^{1}_{S}\bigl(X,\Lie(Y/S,\mathcal{M})\bigr)\).
One has thus obtained the following:

\begin{proposition}{4.2}\label{II.4.2.prop}
Let \(u\colon X\to Y\) be a morphism of \(S\)-groups. One has an isomorphism
of \(S\)-functors, functorial in \(\mathcal{M}\):
\[
L^{u}_{\Hom_{S\text{-gr.}}(X,Y)/S}(\mathcal{M})
\isomto
Z^{1}_{S}\bigl(X,\Lie(Y/S,\mathcal{M})\bigr).
\]
\end{proposition}

{\renewcommand{\thefootnote}{(58)}\nde{One has added the sentences which
follow and proposition~4.2 bis, implicit in the original.}}
Suppose moreover that \(Y/S\) satisfies~(E). Then it follows from~4.2.0.3,
exactly as in the proof of~3.11.1, that
\(\Hom_{S\text{-gr.}}(X,Y)/S\) satisfies~(E). Thus one has (\cf 3.5.1):
\[
L^{u}_{\Hom_{S\text{-gr.}}(X,Y)/S}(\mathcal{M}\oplus\mathcal{N})
\simeq
L^{u}_{\Hom_{S\text{-gr.}}(X,Y)/S}(\mathcal{M})
\times_{S}
L^{u}_{\Hom_{S\text{-gr.}}(X,Y)/S}(\mathcal{N}).
\]
% typo: (Ceci ... a)). -> (Ceci ... a).
(This also follows from~4.2 and~4.2.B~a).)
Consequently,
\(L^{u}_{\Hom_{S\text{-gr.}}(X,Y)/S}(\mathcal{M})\) is equipped, as is
\(Z^{1}_{S}\bigl(X,\Lie(Y/S,\mathcal{M})\bigr)\) (\cf 4.2.B~b)), with an
\(\mathbf{O}_{S}\)-module structure, deduced from the functoriality in
\(\mathcal{M}\). One deduces from this that the isomorphism of~4.2 is, in
this case, an isomorphism of \(\mathbf{O}_{S}\)-modules:

\begin{proposition}{4.2. bis}\label{II.4.2.bis}
\textsuperscript{(58)}\
Let \(u\colon X\to Y\) be a morphism of \(S\)-groups; suppose that \(Y/S\)
satisfies~(E). One has an isomorphism of \(\mathbf{O}_{S}\)-modules,
functorial in \(\mathcal{M}\):
\[
L^{u}_{\Hom_{S\text{-gr.}}(X,Y)/S}(\mathcal{M})
\isomto
Z^{1}_{S}\bigl(X,\Lie(Y/S,\mathcal{M})\bigr).
\]
\end{proposition}

Moreover, when \(Y/S\) satisfies~(E), one deduces from~3.13.1, exactly as in
the proof of~3.13.2, that for every \(u\in\Isom_{S\text{-gr.}}(X,Y)\) one has
an isomorphism functorial in \(\mathcal{M}\):
\[
(*)\qquad
L^{u}_{\Isom_{S\text{-gr.}}(X,Y)/S}(\mathcal{M})
\isomto
L^{u}_{\Hom_{S\text{-gr.}}(X,Y)/S}(\mathcal{M}).
\]
One deduces from this the following two corollaries.

\begin{corollary}{4.2.1}\label{II.4.2.1}
Let \(u\colon X\to Y\) be a morphism of \(S\)-groups; if \(Y/S\)
satisfies~(E), one has an isomorphism of \(\mathbf{O}_{S}\)-modules,
functorial in \(\mathcal{M}\):
\[
L^{u}_{\Isom_{S\text{-gr.}}(X,Y)/S}(\mathcal{M})
\isomto
Z^{1}_{S}\bigl(X,\Lie(Y/S,\mathcal{M})\bigr).
\]
\end{corollary}

\begin{corollary}{4.2.2}\label{II.4.2.2}
Let \(X\) be an \(S\)-group; if \(X/S\) satisfies~(E), one has an isomorphism
of \(\mathbf{O}_{S}\)-modules, functorial in \(\mathcal{M}\):
\[
\Lie\bigl(\Aut_{S\text{-gr.}}(X)/S,\mathcal{M}\bigr)
\isomto
Z^{1}_{S}\bigl(X,\Lie(X/S,\mathcal{M})\bigr).
\]
\end{corollary}

{\renewcommand{\thefootnote}{(59)}\nde{One has added the sentence which
follows.}}
Furthermore, if \(Y\) is \emph{commutative}, the adjoint action of \(Y\) on
\(L=\Lie(X/S,\mathcal{M})\) is trivial, % typo?: kept as printed (adjoint action of Y on Lie(X/S, M))
whence \(Z^{1}_{S}(X,L)=\Hom_{S\text{-gr.}}(X,L)\). Thus:

\begin{corollary}{4.2.3}\label{II.4.2.3}
Let \(Y\) be a commutative \(S\)-group; one has an isomorphism of
\(S\)-functors, functorial in \(\mathcal{M}\):
\[
L^{u}_{\Hom_{S\text{-gr.}}(X,Y)/S}(\mathcal{M})
\isomto
\Hom_{S\text{-gr.}}\bigl(X,\Lie(Y/S,\mathcal{M})\bigr).
\]
\end{corollary}

\numpara{4.3}\label{II.4.3}
Let us now consider the case in which \(X\) and \(Y\) are
\(\mathbf{O}_{S}\)-modules. Let us recall (\cf 4.2.0.4) that one denotes by
\(T'_{Y/S}(\mathcal{M})\) (\resp \(\Lie'(Y/S,\mathcal{M})\)) the functor
\(T_{Y/S}(\mathcal{M})\) (\resp \(\Lie(Y/S,\mathcal{M})\)) equipped with the
\(\mathbf{O}_{S}\)-module structure deduced from that of \(Y\).

When \(Y/S\) satisfies~(E), one will always denote by
\(\Lie(Y/S,\mathcal{M})\) the functor \(\Lie(Y/S,\mathcal{M})\) equipped with
the \(\mathbf{O}_{S}\)-module structure defined for every functor
satisfying~(E). In this case, we know (\cf 3.9) that the abelian group
structures of \(\Lie(Y/S,\mathcal{M})\) and \(\Lie'(Y/S,\mathcal{M})\)
coincide, but the same is not a priori true of the module structures (see a
counter-example in paragraph~6.3). For every \(S'\to S\) and
\(a\in\mathbf{O}(S')\), one will denote by \(a\cdot' m\) (\resp \(a\cdot m\))
the action of \(a\) on \(m\in\Lie'(Y/S,\mathcal{M})(S')\) (\resp on
\(m\in\Lie(Y/S,\mathcal{M})(S')\)), and likewise for the action of \(a\) on
\(T'_{Y/S}(\mathcal{M})\) and \(T_{Y/S}(\mathcal{M})\).

One has \(T'_{Y/S}(\mathcal{M})\simeq\Lie'(Y/S,\mathcal{M})\oplus Y\) as
\(\mathbf{O}_{S}\)-modules; consequently, one obtains, exactly as for
% typo: les proposition -> les propositions
propositions~4.2 and~4.2 bis, the following:

\begin{proposition}{4.3}\label{II.4.3.prop}
Let \(u\colon X\to Y\) be a morphism of \(\mathbf{O}_{S}\)-modules. One has
an isomorphism of \(S\)-functors, functorial in \(\mathcal{M}\):
\[
(*)\qquad
L^{u}_{\Hom_{\mathbf{O}_{S}\text{-mod.}}(X,Y)/S}(\mathcal{M})
\isomto
\Hom_{\mathbf{O}_{S}\text{-mod.}}\bigl(X,\Lie'(Y/S,\mathcal{M})\bigr).
\]
{\renewcommand{\thefootnote}{(60)}\nde{One has added the sentence which
follows.}}
If \(Y/S\) satisfies~(E), then
\(\Hom_{\mathbf{O}_{S}\text{-mod.}}(X,Y)/S\) satisfies~(E) and \((*)\) is an
isomorphism of \(\mathbf{O}_{S}\)-modules when one equips the two members
with the \(\mathbf{O}_{S}\)-module structure deduced from the functoriality
in \(\mathcal{M}\).
{\renewcommand{\thefootnote}{(61)}\nde{Note that on the right-hand term, it
is the structure defined by \((a\varphi)(x)=a\cdot\varphi(x)\), where
\(a\cdot\varphi(x)\) denotes the action of \(a\) on
\(\varphi(x)\in\Lie(Y/S,\mathcal{M})\). This differs from the action
\((a\cdot'\varphi)(x)=a\cdot'\varphi(x)=\varphi(ax)\), but coincides with
the latter if \(\Lie(Y/S,\mathcal{M})=\Lie'(Y/S,\mathcal{M})\).}}
\end{proposition}

\begin{remark}{4.3. bis}\label{II.4.3.bis}
{\renewcommand{\thefootnote}{(62)}\nde{One has added this remark, which will
be useful in~4.5.1.}}
Let \(u\colon X\to Y\) be a morphism of \(\mathbf{O}_{S}\)-modules; let us
denote by \(\tau_{u}\) the map which to every morphism of
\(\mathbf{O}_{S}\)-modules
\(\varphi\colon X\to\Lie'(Y/S,\mathcal{M})\) associates the morphism
\[
u\oplus\varphi\colon X\longrightarrow T'_{Y/S}(\mathcal{M})
=
Y\oplus\Lie'(Y/S,\mathcal{M}).
\]
\end{remark}

Then the isomorphism of~4.3 fits into the following commutative diagram,
functorial in \(\mathcal{M}\):
\[
\xymatrix{
L^{u}_{\Hom_{\mathbf{O}_{S}\text{-mod.}}(X,Y)/S}(\mathcal{M})
  \ar[r]^{\sim}
  \ar@{^{(}->}[d]
&
\Hom_{\mathbf{O}_{S}\text{-mod.}}\bigl(X,\Lie'(Y/S,\mathcal{M})\bigr)
  \ar@{^{(}->}[d]^{\tau_{u}}
\\
T_{\Hom_{\mathbf{O}_{S}\text{-mod.}}(X,Y)/S}(\mathcal{M})
  \ar[r]^{\sim}
&
\Hom_{\mathbf{O}_{S}\text{-mod.}}\bigl(X,T'_{Y/S}(\mathcal{M})\bigr).
}
\]

Moreover, when \(Y/S\) satisfies~(E), one deduces from~3.13.1, exactly as in
the proof of~3.13.2, that for every
\(u\in\Isom_{\mathbf{O}_{S}\text{-mod.}}(X,Y)\), one has
\[
(*)\qquad
L^{u}_{\Isom_{\mathbf{O}_{S}\text{-mod.}}(X,Y)/S}(\mathcal{M})
=
L^{u}_{\Hom_{\mathbf{O}_{S}\text{-mod.}}(X,Y)/S}(\mathcal{M}).
\]

\begin{corollary}{4.3.1}\label{II.4.3.1}
Let \(X\) be an \(\mathbf{O}_{S}\)-module satisfying~(E) with respect to \(S\).
% typo: On un isomorphisme -> On a un isomorphisme
One has a functorial isomorphism in \(\mathcal{M}\):
\[
\Lie\bigl(\Aut_{\mathbf{O}_{S}\text{-mod.}}(X)/S,\mathcal{M}\bigr)
\isomto
\Hom_{\mathbf{O}_{S}\text{-mod.}}\bigl(X,\Lie'(X/S,\mathcal{M})\bigr)
\]
which respects the \(\mathbf{O}_{S}\)-module structures deduced from the
functoriality in \(\mathcal{M}\).
{\renewcommand{\thefootnote}{(63)}\nde{\Cf N.D.E.~(61). On the other hand,
one has added the sentence which follows, as well as its proof.}}
In particular, \(\Aut_{\mathbf{O}_{S}\text{-mod.}}(X)/S\) satisfies~(E).
\end{corollary}

\begin{proof}
The first assertion follows from \((*)\) and~4.3; let us prove the second.
Since \(X/S\) satisfies~(E), one has isomorphisms of
\(\mathbf{O}_{S}\)-modules
\(\Lie'(X/S,\mathcal{M}\oplus\mathcal{N})
\simeq
\Lie'(X/S,\mathcal{M})\times_{S}\Lie'(X/S,\mathcal{N})\), and hence:
\begin{align*}
\Lie\bigl(\Aut_{\mathbf{O}_{S}\text{-mod.}}(X)/S,\mathcal{M}\oplus\mathcal{N}\bigr)
&\simeq
\Lie\bigl(\Aut_{\mathbf{O}_{S}\text{-mod.}}(X)/S,\mathcal{M}\bigr)\\
&\qquad\times_{S}
\Lie\bigl(\Aut_{\mathbf{O}_{S}\text{-mod.}}(X)/S,\mathcal{N}\bigr).
\end{align*}
In view of~4.1.1.2~a), this proves that
\(\Aut_{\mathbf{O}_{S}\text{-mod.}}(X)/S\) satisfies~(E).
\end{proof}

\numpara{4.3.2}\label{II.4.3.2}
Before continuing in this direction, let us examine more closely the
relations between \(Y\), \(\Lie(Y/S)\) and \(\Lie'(Y/S)\). Let us first
remark that
\[
(1)\qquad
\Lie(\mathbf{O}_{S}/S,\mathcal{M})
=
\Lie'(\mathbf{O}_{S}/S,\mathcal{M})
=
W(\mathcal{M})
\]
(where \(W(\mathcal{M})\) is defined in~I~4.6) and that one therefore has a
canonical isomorphism
\[
(2)\qquad
d\colon\mathbf{O}_{S}\isomto\Lie(\mathbf{O}_{S}/S).
\]
Now let \(F\) be an \(\mathbf{O}_{S}\)-module. For every
\(S_{2}\to S_{1}\to S\)
{\renewcommand{\thefootnote}{(64)}\nde{One has changed \(S'\to S\) into
\(S_{2}\to S_{1}\to S\), for one must let \(S_{2}\) and \(S_{1}\) vary
(\cf below). On the other hand, one has expanded the original in what
follows.}}, one has a dihomomorphism
\[
(3)\qquad
\begin{cases}
F(S_{1})\to F(S_{2})\\
\mathbf{O}(S_{1})\to\mathbf{O}(S_{2}),
\end{cases}
\]
whence a morphism of \(\mathbf{O}(S_{2})\)-modules
\[
F(S_{1})\otimes_{\mathbf{O}(S_{1})}\mathbf{O}(S_{2})\longrightarrow F(S_{2}).
\]

In particular, setting \(S_{1}=S'\) and \(S_{2}=I_{S'}(\mathcal{M})\), one
deduces from this morphisms of \(\mathbf{O}(S')\)-modules, functorial in
\(\mathcal{M}\),
\[
F(S')\otimes_{\mathbf{O}(S')}T_{\mathbf{O}_{S}/S}(\mathcal{M})(S')
\longrightarrow
T'_{F/S}(\mathcal{M})(S');
\]
letting \(S'\) vary, one obtains morphisms of \(\mathbf{O}_{S}\)-modules,
functorial in \(\mathcal{M}\),
\[
(4)\qquad
F\otimes_{\mathbf{O}_{S}}T_{\mathbf{O}_{S}/S}(\mathcal{M})
\longrightarrow
T'_{F/S}(\mathcal{M}).
\]
These morphisms are functorial in \(\mathcal{M}\), hence compatible with the
projections of the tangent bundles onto their bases; they therefore define
morphisms of \(\mathbf{O}_{S}\)-modules
\[
(5)\qquad
F\otimes_{\mathbf{O}_{S}}\Lie(\mathbf{O}_{S}/S,\mathcal{M})
\longrightarrow
\Lie'(F/S,\mathcal{M})
\]
such that one has the commutative diagram:
\[
\xymatrix{
0 \ar[r]
&
F\otimes_{\mathbf{O}_{S}}\Lie(\mathbf{O}_{S}/S,\mathcal{M}) \ar[r] \ar[d]
&
F\otimes_{\mathbf{O}_{S}}T_{\mathbf{O}_{S}/S}(\mathcal{M}) \ar[r] \ar[d]
&
F \ar[r]
&
0
\\
0 \ar[r]
&
\Lie'(F/S,\mathcal{M}) \ar[r]
&
T'_{F/S}(\mathcal{M}) \ar[r]
&
F \ar[r]
&
0.
}
\]

One may regard the morphisms~(5) as morphisms of abelian \(S\)-groups
\[
(6)\qquad
F\otimes_{\mathbf{O}_{S}}\Lie(\mathbf{O}_{S}/S,\mathcal{M})
\longrightarrow
\Lie(F/S,\mathcal{M})\,;
\]
by tensoring \(F\) with the isomorphism
\(d\colon\mathbf{O}_{S}\isomto\Lie(\mathbf{O}_{S}/S)\), one deduces from this
(for \(\mathcal{M}=\mathbf{O}_{S}\)) a morphism of abelian \(S\)-groups
\[
(7)\qquad
F\isomto F\otimes_{\mathbf{O}_{S}}\Lie(\mathbf{O}_{S}/S)
\longrightarrow
\Lie(F/S)
\]
also denoted \(d\colon F\to\Lie(F/S)\).

\begin{remark}{4.3.3}\label{II.4.3.3}
{\renewcommand{\thefootnote}{(65)}\nde{One found in the original the
assertion that when \(F/S\) satisfies~(E), the morphisms~(6) (and hence~(7))
are morphisms of \(\mathbf{O}_{S}\)-modules, an assertion contradicted by a
counter-example given in~6.3. One has deleted this assertion, inserted here
the aforementioned example, and modified definition~4.4 accordingly. This is
of no consequence for what follows.}}
When \(F/S\) satisfies~(E), the morphisms~(6) and~(7) are not necessarily
morphisms of \(\mathbf{O}_{S}\)-modules, when one equips the two members
with the module structures deduced from that of \(\mathcal{M}\) by means of
condition~(E).

For example, let \(k\) be a field of characteristic \(p>0\), \(S=\Spec(k)\),
and \(F\) the \(\mathbf{O}_{S}\)-module which to every \(S\)-scheme \(T\)
associates \(F(T)=\Gamma(T,\mathbf{O}_{T})\) equipped with the
\(\mathbf{O}(T)\)-module structure obtained by letting the scalars operate
via the \(p\)-th power, \ie \(r\cdot f=r^{p}f\), for \(r\in\mathbf{O}(T)\),
\(f\in F(T)\). As an \(S\)-group functor, \(F\) is isomorphic to
\(\mathbf{G}_{a,S}\). Thus \(F\) satisfies~(E) and \(\Lie(F/S)\) is identified with
\(\Lie(\mathbf{G}_{a,S}/S)\simeq\mathbf{O}_{S}\). Then the canonical morphism
\(d\colon F\to\Lie(F/S)\) is, for every \(T\to S\), the identity map
\(F(T)\to\mathbf{O}(T)\): it does respect the abelian group structure but
not the \(\mathbf{O}_{S}\)-module structure.
\end{remark}

\begin{remark}{4.3.4}\label{II.4.3.4}
{\renewcommand{\thefootnote}{(66)}\nde{One has added this remark, which will
be useful in~4.6.2.}}
One can make the morphisms~(4) and~(5) explicit as follows. The morphism
\(\Theta\colon F\otimes_{\mathbf{O}_{S}}T_{\mathbf{O}_{S}/S}(\mathcal{M})
\to T'_{F/S}(\mathcal{M})
% typo: Hom_S(I_S(M, F) -> Hom_S(I_S(M), F)
=\Hom_{S}(I_{S}(\mathcal{M}),F)\)
is defined by: for every \(S'\to S\),
\(\alpha\in\mathbf{O}(I_{S'}(\mathcal{M}))\), and \(f\colon S'\to F\),
\[
\Theta(f\otimes\alpha)=\alpha\cdot(\tau_{0}\circ f)=\alpha\cdot(f\circ\rho),
\]
where \(\tau_{0}\) is the zero section \(F\to T'_{F/S}(\mathcal{M})\) and
\(\rho\) the structural morphism \(I_{S'}(\mathcal{M})\to S'\). Then
\(\Theta\) induces a morphism
\(\theta\colon F\otimes_{\mathbf{O}_{S}}\Lie(\mathbf{O}_{S}/S,\mathcal{M})
\to\Lie'(F/S,\mathcal{M})\); this results from the ``functoriality in
\(\mathcal{M}\)'' already mentioned after~(4), and can be seen explicitly as
follows. On the one hand, one has
\[
\Lie'(F/S,\mathcal{M})(S')
=
\bigl\{\varphi\in\Hom_{S}(I_{S'}(\mathcal{M}),F)
\bigm|\varphi\circ\varepsilon_{\mathcal{M}}=e\bigr\},
\]
where \(e\) denotes the unit section \(S'\to F\). On the other hand,
\(\Lie(\mathbf{O}_{S}/S,\mathcal{M})(S')=\Gamma(S',\mathcal{M})\) is the
kernel of the augmentation
\(\eta\colon\mathbf{O}(I_{S'}(\mathcal{M}))\to\mathbf{O}(S')\), and it is
therefore a matter of seeing that if \(f\in F(S')\) and
\(\alpha\in\Gamma(S',\mathcal{M})\), then
\(\alpha\cdot(f\circ\rho)\circ\varepsilon_{\mathcal{M}}=e\).

Let us consider the dihomomorphism~(3), in the case where \(S_{2}\to S_{1}\)
is the \(S\)-morphism
\(\varepsilon_{\mathcal{M}}\colon S'\to I_{S'}(\mathcal{M})\); one then has
a commutative diagram
\[
\xymatrix{
F\bigl(I_{S'}(\mathcal{M})\bigr)
  \ar[r]^{F(\varepsilon_{\mathcal{M}})}
  \ar[d]_{\alpha}
&
F(S')
  \ar[d]^{\eta(\alpha)}
\\
F\bigl(I_{S'}(\mathcal{M})\bigr)
  \ar[r]^{F(\varepsilon_{\mathcal{M}})}
&
F(S')
.}
\]
For every \(\varphi\colon I_{S'}(\mathcal{M})\to F\), one therefore has
\((\alpha\cdot\varphi)\circ\varepsilon_{\mathcal{M}}
=\eta(\alpha)\cdot(\varphi\circ\varepsilon_{\mathcal{M}})\), whence
\((\alpha\cdot\varphi)\circ\varepsilon_{\mathcal{M}}=e\) if
\(\eta(\alpha)=0\).

In particular, taking \(\mathcal{M}=t\,\mathbf{O}_{S}\), one has
\(\Lie(\mathbf{O}_{S}/S)=t\,\mathbf{O}_{S}\) and the morphism
\(F\to\Lie'(F/S)\) is given by \(f\mapsto t\cdot(f\circ\rho)\).
\end{remark}

\begin{remark}{4.3.5}\label{II.4.3.5}
{\renewcommand{\thefootnote}{(67)}\nde{One has added this remark, which will
be useful in~4.5 and~4.7.}}
Let \(F\) be an \(\mathbf{O}_{S}\)-module, set
\(E=\End_{\mathbf{O}_{S}\text{-mod.}}(F)\) and denote by \(d_{F}\) and
\(d_{E}\) the morphisms of \(\mathbf{O}_{S}\)-modules given by~4.3.2~(5):
\begin{align*}
d_{F}&\colon F\otimes_{\mathbf{O}_{S}}\Lie(\mathbf{O}_{S}/S,\mathcal{M})
\longrightarrow\Lie'(F/S,\mathcal{M}),\\
d_{E}&\colon E\otimes_{\mathbf{O}_{S}}\Lie(\mathbf{O}_{S}/S,\mathcal{M})
\longrightarrow\Lie'(E/S,\mathcal{M}).
\end{align*}
One deduces from~4.3.4 the following commutative diagram of morphisms of
\(\mathbf{O}_{S}\)-modules:
\[
\xymatrix@C=4em{
E\otimes_{\mathbf{O}_{S}}\Lie(\mathbf{O}_{S}/S,\mathcal{M})
  \ar[r]^{d_{E}}
&
\Lie'(E/S,\mathcal{M})
\\
\Hom_{\mathbf{O}_{S}}\bigl(F,\,
  F\otimes_{\mathbf{O}_{S}}\Lie(\mathbf{O}_{S}/S,\mathcal{M})\bigr)
  \ar[r]^{d_{F}}
&
\Hom_{\mathbf{O}_{S}}\bigl(F,\Lie'(F/S,\mathcal{M})\bigr)
  \ar[u]_{\simeq\ (*)}
}
\]
where the right-hand vertical arrow is the isomorphism \((*)\) of~4.3.
Thus (\loccit), if \(F/S\) satisfies~(E), then \(E/S\) satisfies~(E) and
\((*)\) is also an isomorphism of \(\mathbf{O}_{S}\)-modules when one equips
the right-hand terms with the \(\mathbf{O}_{S}\)-module structure deduced
from~(E).
\end{remark}

\begin{remark}{4.4.0}\label{II.4.4.0}
{\renewcommand{\thefootnote}{(68)}\nde{In view of the modifications in
definition~4.4 (\cf N.D.E.~(65)), one has inserted here this remark, which
in the original appeared in the proof of~4.4.1.}}
In~4.3.2, the morphisms~(4) are isomorphisms if and only if the
morphisms~(5) are. Moreover, if these conditions are satisfied, then \(F/S\)
satisfies~(E). Indeed, it suffices to check that
\(\Lie'(F/S,\mathcal{M}\oplus\mathcal{N})
\simeq
\Lie'(F/S,\mathcal{M})\times_{S}\Lie'(F/S,\mathcal{N})\).
Now one has the commutative diagram below, in which by hypothesis the
horizontal arrows are isomorphisms:
\[
\xymatrix@C=3em{
F\otimes_{\mathbf{O}_{S}}\Lie(\mathbf{O}_{S}/S,\mathcal{M}\oplus\mathcal{N})
  \ar[r]^{\sim}
  \ar[d]^{\wr}
&
\Lie'(F/S,\mathcal{M}\oplus\mathcal{N})
  \ar[d]
\\
F\otimes_{\mathbf{O}_{S}}
  \bigl(\Lie(\mathbf{O}_{S}/S,\mathcal{M})
    \times_{S}\Lie(\mathbf{O}_{S}/S,\mathcal{N})\bigr)
  \ar[r]^{\sim}
&
\Lie'(F/S,\mathcal{M})\times_{S}\Lie'(F/S,\mathcal{N})\,;
}
\]
the second vertical arrow is therefore also an isomorphism, \ie \(F/S\)
satisfies~(E).
\end{remark}

\begin{definition}{4.4}\label{II.4.4}
One says that \(F\) is a \emph{good} \(\mathbf{O}_{S}\)-module if the
morphisms
\[
F\otimes_{\mathbf{O}_{S}}T_{\mathbf{O}_{S}/S}(\mathcal{M})
\longrightarrow
T_{F/S}(\mathcal{M})
\]
or, equivalently,
\[
F\otimes_{\mathbf{O}_{S}}\Lie(\mathbf{O}_{S}/S,\mathcal{M})
\longrightarrow
\Lie(F/S,\mathcal{M}),
\]
are isomorphisms of abelian \(S\)-groups (so that \(F/S\) satisfies~(E)) and
if moreover they respect the \(\mathbf{O}_{S}\)-module structures deduced
from condition~(E).
\end{definition}

\begin{corollary}{4.4.1}\label{II.4.4.1}
{\renewcommand{\thefootnote}{(69)}\nde{The original showed that~(i) implies
(ii) and~(iii); one has added the implication~(ii)\(\Rightarrow\)(i).}}
Let \(F\) be an \(\mathbf{O}_{S}\)-module. Consider the following conditions:
\begin{enumeratei}
\item \(F\) is a good \(\mathbf{O}_{S}\)-module.
\item \(F/S\) satisfies~(E) and \(d\colon F\to\Lie(F/S)\) is an isomorphism
of \(\mathbf{O}_{S}\)-modules.
\item \(\Lie(F/S,\mathcal{M})=\Lie'(F/S,\mathcal{M})\).
\end{enumeratei}
Then one has \textup{(i)}\(\Leftrightarrow\)\textup{(ii)}\(\Rightarrow\)\textup{(iii)}.
\end{corollary}

\begin{proof}
\textup{(i)}\(\Rightarrow\)\textup{(ii)} follows from the definition. To prove
\textup{(ii)}\(\Rightarrow\)\textup{(i)}, one must show that the morphisms of
abelian \(S\)-groups (functorial in \(\mathcal{M}\))
\[
F\otimes_{\mathbf{O}_{S}}\Lie(\mathbf{O}_{S}/S,\mathcal{M})
\isomto
\Lie(F/S,\mathcal{M})
\]
are isomorphisms of \(\mathbf{O}_{S}\)-modules. Since \(F/S\) satisfies~(E),
the two members transform finite direct sums of copies of \(\mathbf{O}_{S}\)
into finite products of abelian \(S\)-groups. This reduces us to the case
\(\mathcal{M}=\mathbf{O}_{S}\), which follows from the hypothesis.

Finally \textup{(i)}\(\Rightarrow\)\textup{(iii)} follows from the definition
and from the fact that the isomorphisms
\[
F\otimes_{\mathbf{O}_{S}}\Lie(\mathbf{O}_{S}/S,\mathcal{M})
\isomto
\Lie'(F/S,\mathcal{M})
\]
of~4.3.2~(5) are morphisms of \(\mathbf{O}_{S}\)-modules.
\end{proof}

\begin{examples}{4.4.2}\label{II.4.4.2}
For every quasi-coherent \(\mathbf{O}_{S}\)-Module \(\mathcal{E}\), the
\(\mathbf{O}_{S}\)-modules \(V(\mathcal{E})\) and \(W(\mathcal{E})\) defined
in~I~4.6 are good.
\end{examples}

{\renewcommand{\thefootnote}{(70)}\nde{One has added what follows.}}
Indeed, for every \(f\colon S'\to S\), the morphisms
\begin{align*}
V(\mathcal{E})(S')\otimes_{\mathbf{O}(S')}\mathbf{O}\bigl(I_{S'}(\mathcal{M})\bigr)
&\longrightarrow
T_{V(\mathcal{E})/S}(\mathcal{M})(S')\\
W(\mathcal{E})(S')\otimes_{\mathbf{O}(S')}\mathbf{O}\bigl(I_{S'}(\mathcal{M})\bigr)
&\longrightarrow
T_{W(\mathcal{E})/S}(\mathcal{M})(S')
\end{align*}
correspond, respectively, to the morphisms:
\begin{align*}
&\Hom_{\mathbf{O}_{S'}\text{-mod.}}\bigl(f^{*}(\mathcal{E}),\mathbf{O}_{S'}\bigr)
\otimes_{\mathbf{O}(S')}
\Gamma\bigl(S',\mathcal{D}_{\mathbf{O}_{S'}}(\mathcal{M})\bigr)\\
&\qquad\longrightarrow
\Hom_{\mathbf{O}_{S'}\text{-mod.}}\bigl(f^{*}(\mathcal{E}),
  \mathcal{D}_{\mathbf{O}_{S'}}(\mathcal{M})\bigr)\\[0.4em]
&\Gamma\bigl(S',f^{*}(\mathcal{E})\bigr)
\otimes_{\mathbf{O}(S')}
\Gamma\bigl(S',\mathcal{D}_{\mathbf{O}_{S'}}(\mathcal{M})\bigr)\\
&\qquad\longrightarrow
\Gamma\bigl(S',f^{*}(\mathcal{E})
  \otimes_{\mathbf{O}_{S'}}\mathcal{D}_{\mathbf{O}_{S'}}(\mathcal{M})\bigr);
\end{align*}
these are isomorphisms since \(\mathcal{D}_{\mathbf{O}_{S'}}(\mathcal{M})\)
is isomorphic, as an \(\mathbf{O}_{S'}\)-module, to a finite direct sum of
copies of \(\mathbf{O}_{S'}\).

\begin{proposition}{4.5}\label{II.4.5}
Let \(F\) be a good \(\mathbf{O}_{S}\)-module. Then:
\begin{enumeratei}
\item \(\Aut_{\mathbf{O}_{S}\text{-mod.}}(F)/S\) satisfies~(E) and one has a
functorial isomorphism
\[
\Lie\bigl(\Aut_{\mathbf{O}_{S}\text{-mod.}}(F)/S,\mathcal{M}\bigr)
\isomto
\Hom_{\mathbf{O}_{S}\text{-mod.}}\bigl(F,\Lie(F/S,\mathcal{M})\bigr)
\]
which respects the \(\mathbf{O}_{S}\)-module structures deduced from
condition~(E). In particular, one has an isomorphism of
\(\mathbf{O}_{S}\)-modules
\[
\Lie\bigl(\Aut_{\mathbf{O}_{S}\text{-mod.}}(F)/S\bigr)
\isomto
\End_{\mathbf{O}_{S}\text{-mod.}}(F).
\]
\item {\renewcommand{\thefootnote}{(71)}\nde{One has added point~(ii), an
immediate consequence of what precedes.}}
Moreover, \(\End_{\mathbf{O}_{S}\text{-mod.}}(F)\) is a good
\(\mathbf{O}_{S}\)-module.
\end{enumeratei}
\end{proposition}

Indeed, by~4.4.1, \(F/S\) satisfies~(E) and
\begin{equation}
\tag{1}
\Lie(F/S,\mathcal{M})
=
\Lie'(F/S,\mathcal{M})
\simeq
F\otimes_{\mathbf{O}_{S}}\Lie(\mathbf{O}_{S}/S,\mathcal{M}).
\end{equation}
Assertion~(i) then follows from~4.3.1. Set
\(E=\End_{\mathbf{O}_{S}\text{-mod.}}(F)\). By~(1) and remark~4.3.5, one has
the following commutative diagram of morphisms of abelian \(S\)-groups:
\[
\xymatrix@C=3.2em{
\End_{\mathbf{O}_{S}\text{-mod.}}(F)
  \otimes_{\mathbf{O}_{S}}\Lie(\mathbf{O}_{S}/S,\mathcal{M})
  \ar[r]^{d_{E}}
&
\Lie\bigl(\End_{\mathbf{O}_{S}\text{-mod.}}(F)/S,\mathcal{M}\bigr)
\\
\Hom_{\mathbf{O}_{S}}\bigl(F,\,
  F\otimes_{\mathbf{O}_{S}}\Lie(\mathbf{O}_{S}/S,\mathcal{M})\bigr)
  \ar[r]^{d_{F}}
&
\Hom_{\mathbf{O}_{S}}\bigl(F,\Lie(F/S,\mathcal{M})\bigr)
  \ar[u]_{\simeq\ (*)}
}
\]
where \(d_{F}\) and \((*)\) are isomorphisms of \(\mathbf{O}_{S}\)-modules;
consequently, the same is true of \(d_{E}\). This proves~(ii).

\begin{scholium}{4.5.1}\label{II.4.5.1}
{\renewcommand{\thefootnote}{(72)}\nde{One has added this scholium, implicit
in the original, which will be useful in~4.7.1. (Here and in what follows,
one denotes by \(t\), \(t'\), etc., variables of square zero, the letter
\(\varepsilon\) being reserved for the unit section of groups.)}}
Set (\cf 2.1) \(\mathbf{O}_{I_{S}}=\mathbf{O}_{S}\oplus t\,\mathbf{O}_{S}\)
(with \(t^{2}=0\)), and let \(F\) be a good \(\mathbf{O}_{S}\)-module. Then,
for every \(S'\to S\), the morphism
\[
F(S')\oplus t\,F(S')
=
F(S')\otimes_{\mathbf{O}(S')}\mathbf{O}(I_{S'})
\longrightarrow
F(I_{S'})
=
F(S')\oplus\Lie(F/S)(S')
\]
(which is the identity on \(F(S')\)) induces an isomorphism of
\(\mathbf{O}(S')\)-modules
\(t\,F(S')\simeq\Lie(F/S)(S')\).
\end{scholium}
